% Standalone manuscript. Compile twice with pdflatex, or once with tectonic.
\documentclass[11pt,letterpaper]{article}
\usepackage[margin=1in]{geometry}
\usepackage{amsmath,amssymb,amsthm,booktabs}
\usepackage[hidelinks]{hyperref}
\hypersetup{pdftitle={Local error bounds for consecutive NOMAD Taylor steps},
 pdfsubject={A companion to John Atkinson's polynomial ODE derivation}}
\newcommand{\R}{\mathbb R}
\newcommand{\F}{\mathcal D_F}
\newcommand{\file}[1]{\texttt{\detokenize{#1}}}
\newcommand{\abs}[1]{\left|#1\right|}
\newcommand{\norm}[1]{\left\lVert#1\right\rVert}
\newtheorem{theorem}{Theorem}
\newtheorem{proposition}[theorem]{Proposition}
\setlength{\emergencystretch}{1.5em}
\raggedbottom
\begin{document}
\begin{center}
{\LARGE\bfseries Local error bounds for consecutive\\[.12em] NOMAD Taylor steps\par}
\vspace{.65em}
{\large Companion to John Atkinson's 2016--2017 NOMAD notebooks\par}
\vspace{.2em}
{\small Working draft for author review\par}
\end{center}
\begin{abstract}
Two consecutive NOMAD steps, taken from the same state, differ by one
explicitly computable Taylor term. We enclose the error of the lower-order
step in a signed interval centered at that difference. The radius comes from
one further normalized derivative polynomial on a box containing the
solution. A strict inequality on the field coefficients verifies the box.
Both endpoints can be computed rationally for a polynomial system with
binary floating-point data, including the correction for rounding in the
returned solver values. Lean checks the derivation and a concrete numerical
enclosure. An example with a vanishing increment shows why the additional
derivative is necessary for this guarantee.
\end{abstract}

\section{The error seen by two consecutive orders}
Let \(x'=F(x)\) be an autonomous polynomial system in \(m\) real variables,
with exact starting value \(x(a)=z\). Consider a forward step \(h>0\).
Write \(H_{0,\ell}(y)=y_\ell\), \(H_{n+1,\ell}=\F H_{n,\ell}\), where
\(\F P=\sum_j F_j\,\partial_j P\). The NOMAD construction~\cite{notebook}
computes the coefficients and exponents of these polynomials. Its normalized
tables and its degree-\(p\) step satisfy
\begin{equation}\label{eq:step}
 C_{n,\ell}=H_{n,\ell}/n!,\qquad C_{n+1,\ell}=\frac{\F C_{n,\ell}}{n+1},
 \qquad T_p(z,h)=\sum_{n=0}^{p}h^n C_n(z).
\end{equation}
All terms use the same \(z\). Define the signed local error
\(e_{p,\ell}=x_\ell(a+h)-T_{p,\ell}(z,h)\) and the observable increment
\begin{equation}\label{eq:delta}
 d_{p,\ell}=T_{p+1,\ell}(z,h)-T_{p,\ell}(z,h)
             =h^{p+1}C_{p+1,\ell}(z).
\end{equation}

\begin{theorem}\label{th:main}
Suppose a solution exists on \([a,a+h]\) and
\(\abs{C_{p+2,\ell}(x(t))}\le B_{p+2,\ell}\) throughout that interval.
Set \(R_\ell=h^{p+2}B_{p+2,\ell}\). Then
\begin{align}
  |e_{p+1,\ell}| &\le R_\ell, &
  d_{p,\ell}-R_\ell &\le e_{p,\ell}\le d_{p,\ell}+R_\ell,\label{eq:signed}\\
  \max\{0,|d_{p,\ell}|-R_\ell\}
     &\le |e_{p,\ell}| \le |d_{p,\ell}|+R_\ell. &&\label{eq:absolute}
\end{align}
\end{theorem}
\begin{proof}
Along a solution, the chain rule gives
\(\frac{d}{dt}P(x(t))=(\F P)(x(t))\). Induction therefore gives
\(x^{(n)}_\ell(t)=n!C_{n,\ell}(x(t))\). Scalar Taylor's theorem at order
\(p+1\) supplies a \(\xi_\ell\in(a,a+h)\) such that
\(e_{p+1,\ell}=h^{p+2} C_{p+2,\ell}(x(\xi_\ell))\).
Each coordinate can have its own \(\xi_\ell\).
Now \(e_{p,\ell}=d_{p,\ell}+e_{p+1,\ell}\). The triangle and reverse
triangle inequalities give \eqref{eq:absolute}.
\end{proof}

For a component tolerance \(\epsilon_\ell\), the upper endpoint in
\eqref{eq:absolute} certifies that the step meets tolerance whenever it is at
most \(\epsilon_\ell\). A lower endpoint above \(\epsilon_\ell\) certifies a
miss. When the interval straddles the tolerance, an additional order or a
smaller step can resolve the comparison.

\clearpage
\section{Obtaining the radius from NOMAD tables}
For radii \(r_j>0\), define the closed box
\(Q=\{y: |y_j-z_j|\le r_j\text{ for all }j\}\) and put
\(\rho_j=|z_j|+r_j\). For a scalar polynomial
\(P(y)=\sum_u c_u y^u\), define
\begin{equation}\label{eq:maj}
  \mathcal B(P,\rho)=\sum_u |c_u|\prod_j\rho_j^{u_j}.
\end{equation}
Use the convention \(0^0=1\) for an absent factor. Each monomial is
bounded by the corresponding nonnegative product, so
\(\abs{P(y)}\le\mathcal B(P,\rho)\) for \(y\in Q\).
Collecting equal exponents before taking absolute values preserves
cancellations and can reduce the radius.

\begin{proposition}\label{prop:box}
Set \(M_j=\mathcal B(F_j,\rho)\). If
\begin{equation}\label{eq:box}
                  h M_j < r_j\qquad (1\le j\le m),
\end{equation}
the IVP has a unique solution throughout the step and that solution stays
in the box. In Theorem~\ref{th:main} one may use
\(B_{p+2,\ell}=\mathcal B(C_{p+2,\ell},\rho)\).
\end{proposition}
\begin{proof}
A polynomial field is locally Lipschitz, so a unique maximal solution starts
at \(z\). If it first touches a face at time \(u\le a+h\), then before
that time its speed in coordinate \(j\) is at most \(M_j\).
Integration gives \(\abs{x_j(u)-z_j}\le(u-a)M_j\le hM_j<r_j\)
for every \(j\), contradicting contact. The same estimate holds up to any
earlier end of the maximal interval. On the compact box the speed is
bounded, so the solution then has a limit there and extends by local
existence. This gives existence through the endpoint. Applying \eqref{eq:maj}
to the normalized derivative polynomial now gives the asserted value of
\(B_{p+2,\ell}\). See also the local existence and continuation theorems
in Teschl~\cite{teschl}.
\end{proof}

The procedure requires the field table and two more orders beyond \(p\).
The table at order \(p+1\) is evaluated at \(z\) with its signs. At order
\(p+2\) one sums the absolute coefficients of each row against \(\rho^u\).
Both operations use the same exponent columns as the integrator.
The box condition is sufficient. A failure of \eqref{eq:box} leaves this
step uncertified. Taking a shorter step with the same radii reduces both
\(hM_j\) and the remainder radius.

\paragraph{An increment can hide the entire error.}
For any \(p\ge1\), take \(u'=1,\ v'=u^{p+1}\) with \((u,v)(0)=(0,0)\).
Then \(u(h)=h\) and \(v(h)=h^{p+2}/(p+2)\). Both orders return \((h,0)\).
Their difference is zero in every coordinate although the exact error in
\(v\) is positive. Here \(C_{p+2,v}=1/(p+2)\) everywhere. The radius in
\eqref{eq:signed} gives the exact magnitude of the missed term. Conversely,
when \(C_{p+1,\ell}(z)\ne0\), \eqref{eq:signed} gives
\(e_{p,\ell}/d_{p,\ell}=1+O(h)\) as \(h\downarrow0\) in a fixed box.

\section{More orders and floating-point endpoints}
For \(q\ge p\), accumulate the block
\(D_{p,q}=\sum_{n=p+1}^{q} h^n C_n(z)\), with the empty sum equal to zero.
Since \(e_{p,\ell}-D_{p,q,\ell}=e_{q,\ell}\), the same proof yields
\begin{equation}\label{eq:block}
 \abs{e_{p,\ell}-D_{p,q,\ell}}
   \le h^{q+1}\mathcal B(C_{q+1,\ell},\rho).
\end{equation}
Each extra derivative can move the center of the interval as well as change
its radius. Every valid interval encloses the same signed lower-order error,
so their intersection is valid. The componentwise bounds also give
\(\norm{e_p}_\infty\le\max_\ell(|d_{p,\ell}|+R_\ell)\).

\clearpage
\paragraph{The returned PyTorch value.}
Write \(Y_p,Y_{p+1}\) for the actual floating-point results of two calls to the
solver from the same represented initial state. Interpret those finite
outputs as exact binary rationals. Let
\(\widehat d=Y_{p+1}-Y_p\), taking the difference in exact arithmetic, and
let \(s=T_{p+1}(z,h)-Y_{p+1}\) be the signed correction computed by a rational
evaluation of the polynomial tables. Coordinatewise,
\begin{equation}\label{eq:rounded}
  \boxed{\quad x(a+h)-Y_p\ \in\ \widehat d+s+[-R,R],\qquad
          x(a+h)-Y_{p+1}\ \in\ s+[-R,R].\quad}
\end{equation}
Here \([-R,R]\) denotes a componentwise interval. Indeed, subtracting
\(\widehat d+s\) from the low-order error gives exactly \(x(a+h)-T_{p+1}\).
The correction also accounts for rounding that occurs while constructing the
tables. As a simple test, take \(x'=1, z=1, h=2^{-54}\). In binary64 both
orders round to \(1\). The Taylor remainder is zero and the correction is
\(2^{-54}\). The signed interval for each actual output is the singleton
\(\{2^{-54}\}\).

\paragraph{Local and accumulated error.}
Each enclosure starts an exact trajectory at the supplied state. At a later
numerical step this state can differ from the original IVP trajectory.
For example, if the flow from the two initial states has Lipschitz bound
\(L_n\) over step \(n\), an existing global error bound \(E_n\) and a local
endpoint bound \(\tau_n\) give \(E_{n+1}\le L_n E_n+\tau_n\) by the triangle
inequality. The certificates below compute \(\tau_n\). A user certifying
a full trajectory must also enclose both flows and supply their stability
bounds.

\section{An interval that can be checked by hand}
Take \(x'=x^2\), \(z=1\), \(h=1/8\), and orders two and three.
Here \(C_n(y)=y^{n+1}\) for every \(n\ge0\), since
\(y^2\partial_y(y^{n+1})/(n+1)=y^{n+2}\).
With \(r=1/2\) one gets \(\rho=3/2\), \(M=9/4\), and \(hM=9/32<1/2\).
The outputs and interval are
\begin{align*}
  Y_2=T_2&=73/64, &Y_3=T_3&=585/512, &d_2&=1/512,\\
  B_4&=(3/2)^5, &R&=243/131072, &s&=0.
\end{align*}
\begin{equation}\label{eq:example}
       \frac{13}{131072}\ \le\ x(1/8)-Y_2\ \le\ \frac{499}{131072}.
\end{equation}
The exact solution \(x(h)=1/(1-h)\) gives error \(1/448\), which lies in this
interval. The lower endpoint is positive, so the certificate also determines
the sign of the error. Lean derives \eqref{eq:example} directly from the
ODE and the starting state. The companion program obtains the same endpoints
from actual calls to \file{NomadSolver.step}.

\paragraph{The notebook's two-coordinate example.}
For \(F(x,y)=(y^2-x,x^2-y)\) at \((1,2)\), use \(h=1/32\) and radii
\((1/2,1/2)\). The speed bounds are \(31/4\) and \(19/4\), so both box
tests pass. For orders three and four the exact increment is
\(h^4(-91,-49)/24\), obtained from the fourth derivative in the notebook
whitepaper. The test suite checks the rational order-four result in both
coordinates. It also exercises duplicate input columns, shrinking derivative
tables, opposite error signs, a zero field, and rejection at equality in
\eqref{eq:box}.

\clearpage
\section{Reproducing the bounds}
The accompanying Lean files and their mathematical roles are listed below.
The project pins Lean~4 and mathlib to v4.19.0. The proof uses mathlib's
scalar Taylor theorem with Lagrange remainder~\cite{mathlib}. Within-interval
derivatives give the appropriate derivatives at both endpoints.
\begin{center}\small
\begin{tabular}{@{}lp{.53\linewidth}@{}}
\toprule
Declaration in namespace \file{Nomad} & Statement checked\\
\midrule
\file{field_paths}, \file{jet_derivative} & Expanded NOMAD polynomial and its time derivative interpretation.\\
\file{normalized_succ}, \file{step_eq_taylor} & Factorial-scaled tables and the Taylor polynomial.\\
\file{local_remainder}, \file{block_enclosure} & Scalar remainders for one order and for a block.\\
\file{table_bound}, \file{stays_in_box} & Polynomial majorant and first-exit estimate.\\
\file{boxed_consecutive} & Signed local enclosure from both tables and the box test.\\
\file{absolute_enclosure}, \file{rounded_enclosure} & Absolute errors and transfer to the returned values.\\
\file{Example.quadratic_certificate} & Both rational endpoints in \eqref{eq:example}.\\
\bottomrule
\end{tabular}
\end{center}
The interval theorems in Lean quantify over a solution of the ODE on the
step. Proposition~\ref{prop:box} supplies existence and continuation in the
paper. The CPU solver supplies both numerical endpoints. The companion
module \file{nomad_error_bounds.py} reconstructs the normalized tables with
Python's \file{Fraction}. The general Python program and PyTorch kernels are
outside the Lean formalization. Exact fraction arithmetic in the certificate
and tests applies to the field, initial state and step as represented on input.

With Lean's toolchain manager and Python with PyTorch installed, run the
following commands from the repository root. The example and tests use CPU
and were checked with PyTorch 2.6.0.
\begin{verbatim}
(cd lean && lake exe cache get && lake build)
python3 pytorch/nomad_error_bounds.py
python3 -m unittest -b -v pytorch.test_nomad_error_bounds
\end{verbatim}
For other fields, call \file{certify_step(a, k, z, h, p, radii)}. Rows of
\(a,k\) index the state coordinates and columns index the monomials. Use finite
binary64 tensors for \(a,z\) and the radii, nonnegative integer exponents,
a positive floating-point \(h\), and integer \(p\ge1\). The optional keyword
\file{high_order=q} uses the longer block in \eqref{eq:block} for \(q>p\). Each returned record
gives the signed low- and high-order intervals and the absolute low-order
interval. The independent rational table construction is intended for small
systems and modest orders, where its cost is practical.

\begin{thebibliography}{9}\raggedright
\bibitem{notebook} J. Atkinson.
\emph{NOMAD: A Taylor solver for multivariate polynomial differential
equations}. Companion edition of the 2016--2017 notebooks, \file{NOMAD.tex}.
\bibitem{teschl} G. Teschl. \emph{Ordinary Differential Equations and
Dynamical Systems}. American Mathematical Society, 2012, §§2.2 and 2.6.
\href{https://www.mat.univie.ac.at/~gerald/ftp/book-ode/ode.pdf}{Author's edition}.
\bibitem{mathlib} The mathlib community. \emph{Taylor's theorem}, mathlib4
v4.19.0, declaration \file{taylor_mean_remainder_lagrange}.
\href{https://github.com/leanprover-community/mathlib4/blob/v4.19.0/Mathlib/Analysis/Calculus/Taylor.lean}{Versioned source}.
\end{thebibliography}
\end{document}
